% Λύση του 43ου προβλήματος της ενότητας «Άλυτα Προβλήματα».
\subsection*{Λύση Προβλήματος 43 --- Κατακόρυφη παραβολική καμπύλη}

\begin{alist}
  \item Έχουμε
        \[
          g(x)=ax^2+\beta x+\gamma.
        \]
        Άρα
        \[
          g'(x)=2ax+\beta.
        \]

        Από τη συνθήκη $g(0)=100$ παίρνουμε
        \[
          \gamma=100.
        \]
        Από τη συνθήκη $g'(0)=0{,}04$ παίρνουμε
        \[
          \beta=0{,}04.
        \]
        Τέλος, από τη συνθήκη $g'(200)=-0{,}02$ έχουμε
        \[
          2a\cdot200+\beta=-0{,}02.
        \]
        Επειδή $\beta=0{,}04$, προκύπτει
        \[
          400a+0{,}04=-0{,}02.
        \]
        Άρα
        \[
          400a=-0{,}06
          \quad\Rightarrow\quad
          a=-0{,}00015.
        \]

        Επομένως
        \[
          {
          a=-0{,}00015,\qquad
          \beta=0{,}04,\qquad
          \gamma=100
          }.
        \]
        Η εξίσωση της κατακόρυφης καμπύλης είναι
        \[
          {g(x)=-0{,}00015x^2+0{,}04x+100}.
        \]

  \item Το υψόμετρο στο τέλος της καμπύλης είναι
        \[
          g(200)
          =
          -0{,}00015\cdot200^2+0{,}04\cdot200+100.
        \]
        Υπολογίζουμε:
        \[
          g(200)
          =
          -0{,}00015\cdot40000+8+100
          =
          -6+108
          =
          102.
        \]
        Άρα
        \[
          {g(200)=102\si{m}}.
        \]

  \item Η αρχική ευθεία κλίσης περνά από το $A(0,100)$ και έχει κλίση
        $0{,}04$. Επομένως η εξίσωσή της είναι
        \[
          y=100+0{,}04x.
        \]
        Η παραβολική καμπύλη στο $x=0$ έχει
        \[
          g(0)=100
        \]
        και
        \[
          g'(0)=0{,}04.
        \]
        Δηλαδή περνά από το ίδιο σημείο και έχει την ίδια κλίση με την
        αρχική ευθεία. Άρα οι δύο γραμμές έχουν κοινή εφαπτομένη στο
        $A$ και η σύνδεση είναι ομαλή ως προς την κλίση.
\end{alist}
